\section{Avaliação Topológica}
\label{sec:avaliacao}

A avaliação topológica segue um pipeline de nove fases implementado em
\texttt{scripts/voronoi\_foliation\_drm.py}. Primeiro, são extraídas
coordenadas DRM, diagonal efetiva de $G(x)$, gamma e massa por token. Em
seguida, aplica-se KMeans para formar uma tesselação de Voronoi no manifold.
Cada célula é analisada por PCA local, produzindo uma estimativa de dimensão
efetiva por LTSA.

A coerência tangente é calculada por ângulos principais entre subespaços
locais de células vizinhas. A análise de Reeb usa funções escalares como gamma
ou coordenadas do manifold para detectar divisões e fusões de níveis. A
homologia persistente é computada por \emph{ripser} em subsamples de tamanho
controlado, usando normalização z-score e uma regra de barras longas baseada em
fração da maior persistência:
\begin{equation}
  \operatorname{long}(b) \iff
  \operatorname{pers}(b) \geq \rho \max_b \operatorname{pers}(b).
\end{equation}

Para evitar validação por acaso, a homologia é repetida em múltiplos
subsamples. A fração estável de toro é definida por
\begin{equation}
  f_{T^2} =
  \frac{1}{R}\sum_{r=1}^{R}
  \mathbf{1}[H_1^{(r)}=2 \land H_2^{(r)}=1].
\end{equation}
Neste trabalho, a assinatura toroidal é considerada estável quando
$f_{T^2}\geq 0{,}60$.

Também são avaliados modelos nulos: coordenadas embaralhadas por eixo e
coordenadas uniformes. Esses controles ajudam a distinguir estrutura
persistente de artefatos de clustering ou de distribuição marginal.
